\scanpage{64}
On peut montrer directement que $\sigma z=e^{\int\zeta z\,dz}$ est uniforme, \deleted{est donc} \inserted{est même} une fonction entière ; \deleted{il est} on peut aussi procéder directement.

On a \deleted{en effet} $\displaystyle\unsure{p|\wp}z=\frac1{z^2}+\sum{\sum}'\left[\frac1{(z-\omega)^2}-\frac1{\omega^2}\right]$ la \unsure{somme|série} étant unift convergente pour $z$ borné, à distance $>\varepsilon$ des $\omega$. Les dérivées successives \unsure{s'obtient|s'obtiennent} par dérivation terme à terme, la convergence est comme précédemment.

Notons que l'on obtient ainsi les coefficients du développement de $pz$
\[
c_1=3\sum{\sum}'\frac1{\omega^{2}}\qquad c_2=5\sum{\sum}'\frac1{\omega^6}\qquad\mathord{---}
\]

On a
\[
\zeta z=\frac1z+\sum{\sum}'\left[\frac1{z-\omega}+\frac1\omega+\frac{z}{\deleted{\unsure{2}}\omega^2}\right]
\]

Posons enfin
\[
\sigma z=z\prod{\prod}'\left(1-\frac z\omega\right)e^{\frac z\omega+\frac{z^2}{2\omega^2}}
\]

Le produit converge \deleted{et} uniformément \inserted{et absol\textsuperscript{t}} si $z$ est borné ; $\sigma z$ est une fonction entière, qui s'annule pour les $z=\omega$ et pas d'autres points. Elle a été construite de sorte à avoir
\[
\frac{\sigma'z}{\sigma z}=\zeta z
\]

(À remarquer que, si on voulait construire à priori une fonction entière s'annulant aux points $\omega$, aussi simple que possible, c'est $\sigma z$ qu'on aurait trouvé)

$\sigma z$ est une fonction paire. On a
\[
\sigma(z+2\omega_1)=-e^{2\eta_1(z+\omega_1)}\sigma z\qquad
\sigma(z+\omega_2)=-e^{2\eta_2(z+\omega_{\unsure{2|1}})}\sigma z
\]
de $\unsure{\overline{\mathrm m}|=}$ pour $\omega_3$

% archive sheet number: 63
